% Fragmento maestro: incluir desde el anfitrion, con \HMTFGRoot fijada a la raiz del paquete.
% Requiere amsmath, amssymb/mathrsfs o unicode-math, amsthm, booktabs, longtable, hyperref, needspace.
% El anfitrion conserva sus entornos theorem, proposition, lemma y proof.
\providecommand{\HMTFGRoot}{.}
\begingroup
\providecommand{\ZZ}{}\renewcommand{\ZZ}{\mathbb Z}
\providecommand{\RR}{}\renewcommand{\RR}{\mathbb R}
\providecommand{\FF}{}\renewcommand{\FF}{\mathbb F}
\providecommand{\APP}{}\renewcommand{\APP}{\mathrm{APP}}
\providecommand{\TPK}{}\renewcommand{\TPK}{\mathrm{TPK}}
\providecommand{\TRIT}{}\renewcommand{\TRIT}{\{-1,0,+1\}}
% unicode-math usa lBrack/rBrack; newtxmath ya proporciona llbracket/rrbracket.
\providecommand{\llbracket}{\lBrack}
\providecommand{\rrbracket}{\rBrack}
\input{\HMTFGRoot/latex/00_antecedentes_operatorios.tex}
\input{\HMTFGRoot/latex/01_perfil_y_retornos.tex}
\input{\HMTFGRoot/latex/02_profundidad_y_memoria.tex}
\input{\HMTFGRoot/latex/10_lectores_minimos_y_restitucion.tex}
\input{\HMTFGRoot/latex/03_transporte_ordenado_y_jacobi.tex}
\input{\HMTFGRoot/latex/04_clasificacion_y_selector.tex}
\input{\HMTFGRoot/latex/05_entrada_transversal_y_escalado_local.tex}
\input{\HMTFGRoot/latex/06_aislamiento_y_escalado_global.tex}
\input{\HMTFGRoot/latex/07_identificacion_analitica_de_centros.tex}
\input{\HMTFGRoot/latex/08_interpretacion_estructural.tex}
\endgroup
